\documentclass{article}
\usepackage[colorlinks]{hyperref}
% Use the following line for the initial blind version submitted for review:
\usepackage{neurips_2026}

% to compile a preprint version, e.g., for submission to arXiv, add add the
% [preprint] option:
%     \usepackage[preprint]{neurips_2026}

% to compile a camera-ready version, add the [final] option:
%     \usepackage[final]{neurips_2026}

% to avoid loading the natbib package, add option nonatbib:
%    \usepackage[nonatbib]{neurips_2026}

% --- Essential Packages ---
\usepackage{times}
\usepackage{soul}
\usepackage{url}
\usepackage[utf8]{inputenc}
\usepackage{caption}
\usepackage{graphicx}
\usepackage{amsmath}
\usepackage{amsthm}
\usepackage{amssymb}
\usepackage{booktabs}
\usepackage{algorithm}
\usepackage{algorithmic}
\usepackage{bm}
\usepackage{mathtools}
\usepackage{enumitem}
\usepackage{xcolor}

% --- Custom Math Definitions for Physics-AI ---
\newtheorem{theorem}{Theorem}
\newtheorem{definition}{Definition}
\newtheorem{proposition}{Proposition}

\newcommand{\RSMI}{\mathcal{I}_{\text{RSMI}}}
\newcommand{\fermigap}[1]{\Delta_{\text{F},#1}}
\newcommand{\betafunc}{\beta(g)}
\newcommand{\ket}[1]{|#1\rangle}
\newcommand{\bra}[1]{\langle#1|}
% 显示为红色占位，不再出现“TODO”字样
\newcommand{\TODO}[1]{\textcolor{red}{#1}}

\urlstyle{same}

\title{Fermi-Transformers: Physics-Inspired Gapped Attention \\ for Noise-Robust Reasoning in Language Models}

% The \author macro works with any number of authors. There are two commands
% used to separate the names and addresses of multiple authors: \And and \AND.
%
% Using \And between authors leaves it to LaTeX to determine where to break the
% lines. Using \AND forces a line break at that point. So, if LaTeX puts 3 of 4
% authors names on the first line, and the last on the second line, try using
% \AND instead of \And before the third author name.

\author{%
  Anonymous Authors \\
  Anonymous Institution \\
  \texttt{contact@anonymous.edu} \\
}

% Fix for date suppression
\date{}

\begin{document}

\maketitle

% =========================================================
% Abstract
% =========================================================
\begin{abstract}
Large Language Models (LLMs) are notoriously brittle: small syntactic perturbations in the input can trigger disproportionate changes in the output, leading to logical inconsistencies and hallucinations. We view this phenomenon as an instance of \emph{ultraviolet (UV) instability}, where microscopic token-level noise percolates through all-pass attention layers and melts the macroscopic logical structure. 

Inspired by the Renormalization Group (RG) and Fermi statistics in condensed matter physics, we propose the \emph{Fermi-Transformer}, a physics-inspired architecture that selectively integrates out UV noise while preserving long-range semantic structure. Concretely, we introduce (i) an \emph{RSMI-Pooling} module that learns coarse-grained representations by maximizing real-space mutual information between neighboring token blocks, and (ii) a \emph{Fermi-Dirac Gapped Attention} kernel that imposes a learnable energy gap on attention weights, suppressing low-information interactions.

We evaluate Fermi-Transformers on synthetic and real-world reasoning benchmarks with adversarial token-level noise. Across multiple tasks, our model reduces a formalized \emph{logic drift} metric by up to \TODO{xx\%} compared to a vanilla Transformer with comparable capacity, while maintaining or improving clean-task accuracy. Spectral analysis further reveals a stable power-law decay of attention singular values, indicative of a scale-invariant fixed point. Our results suggest that RG and topological principles provide a promising route towards more noise-robust and reliable language models.
\end{abstract}

% =========================================================
% 1. Introduction
% =========================================================
\section{Introduction}

Scaling up Transformer-based LLMs has led to emergent behaviors such as in-context learning and non-trivial reasoning capabilities~\cite{kaplan2020scaling}. However, these systems remain surprisingly fragile: small, seemingly irrelevant perturbations in the prompt can dramatically alter the final answer or even induce hallucinations. This mismatch between local token perturbations and global logical behavior is reminiscent of phenomena in many-body physics, where microscopic fluctuations can destabilize macroscopic phases in the absence of an energy gap.

Following Anderson's perspective that ``more is different''~\cite{anderson1972more}, we argue that robust reasoning should be understood as an emergent \emph{phase} of information flow, governed by effective laws distinct from token-level statistics. From this viewpoint, current Transformers suffer from an \emph{ultraviolet (UV) instability}: standard softmax attention acts as an all-pass filter, allowing high-frequency syntactic noise to propagate across layers and eventually disrupt the infrared (IR) logical structure.

\subsection{Hallucinations as UV Instability}

We formalize this intuition by distinguishing between microscopic token perturbations and macroscopic task behavior. Let $x_{1:T}$ denote a token sequence and $f_\theta$ a language model. A \emph{UV perturbation} is generated by a stochastic or adversarial corruption operator $\mathcal{C}$ acting on a small subset of positions:
\begin{equation}
    \tilde{x}_{1:T} = \mathcal{C}(x_{1:T}), 
    \quad
    \big|\{t : \tilde{x}_t \neq x_t\}\big| \ll T.
\end{equation}
The downstream task metric (e.g., the final numeric answer in a math word problem, or the program behavior in code generation) characterizes the IR logical structure.

We define the \emph{logic drift} of $f_\theta$ under perturbation distribution $\mathcal{C}$ as
\begin{equation}
    \mathrm{Drift}(f_\theta) =
    \mathbb{E}_{x,\mathcal{C}}\big[
        d\big(f_\theta(x), f_\theta(\mathcal{C}(x))\big)
    \big],
    \label{eq:logic-drift}
\end{equation}
where $d(\cdot,\cdot)$ measures discrepancy between outputs (e.g., task loss, edit distance, or semantic divergence). Empirically, we observe that standard Transformers can exhibit large $\mathrm{Drift}(f_\theta)$ even when performance on clean inputs is high, indicating a UV-instability regime.

\subsection{Physics-Inspired Approach}

The Renormalization Group (RG) provides a successful framework for understanding how microscopic fluctuations affect macroscopic phases~\cite{wilson1975renormalization}. Recent work has started to connect RG ideas with deep learning~\cite{mehta2014exact,koch2018mutual,hu2024learning}, and with the geometry and topology of neural representations~\cite{battiloro2024topological,fujimoto2024topological}. Motivated by these developments, we ask:

\begin{quote}
\emph{Can we endow Transformers with an explicit ``energy gap'' in attention, so that UV syntactic noise cannot easily melt the IR logical phase?}
\end{quote}

To this end, we propose the \textbf{Fermi-Transformer}, which incorporates two key components:

\begin{enumerate}[leftmargin=*,topsep=0pt]
    \item \textbf{RSMI-Pooling.} A coarse-graining module that learns block-level representations by maximizing real-space mutual information (RSMI) between neighboring regions~\cite{koch2018mutual}, systematically integrating out locally unpredictable noise.
    \item \textbf{Fermi-Dirac Gapped Attention.} A generalized attention kernel based on Fermi-Dirac statistics, introducing a learnable gap that sharply separates high-information from low-information interactions, thereby stabilizing the semantic topology of attention.
\end{enumerate}

\subsection{Contributions}

This paper makes the following contributions:
\begin{itemize}[leftmargin=*]
    \item We propose a physics-inspired viewpoint that frames LLM hallucinations as \emph{UV instabilities} in semantic flow, and we formalize a \emph{logic drift} metric that quantitatively captures the sensitivity of model outputs to local token perturbations.
    \item We introduce the \emph{Fermi-Transformer}, comprising (i) an RSMI-Pooling layer that learns coarse-grained token blocks by maximizing real-space mutual information, and (ii) a Fermi-Dirac Gapped Attention kernel that imposes a learnable energy gap on attention weights, suppressing low-information interactions while preserving strong logical links.
    \item We demonstrate on noisy reasoning benchmarks that Fermi-Transformers consistently reduce logic drift by up to \TODO{xx\%} relative to strong Transformer baselines at comparable parameter budgets, while maintaining competitive performance on clean inputs. Spectral analysis of the attention matrices further reveals a robust power-law decay of singular values, suggesting convergence towards a scale-invariant fixed point.
\end{itemize}

Together, these results highlight the potential of RG and topological principles as design tools for robust and scalable language models.

% =========================================================
% 2. Background and Problem Formulation
% =========================================================
\section{Background and Problem Formulation}

\subsection{Transformers and All-Pass Attention}

Given a sequence of token embeddings $X \in \mathbb{R}^{T \times d}$, a standard multi-head attention layer computes
\begin{equation}
    \mathrm{Attn}(Q,K,V) = \mathrm{softmax}\!\left(\frac{QK^\top}{\sqrt{d_k}}\right) V,
\end{equation}
where $Q = XW^Q$, $K = XW^K$, $V = XW^V$. The softmax kernel produces strictly positive weights over all positions, effectively implementing an \emph{all-pass} mixing of contextual information. While this is expressive, it also means that small perturbations in $QK^\top$ can globally reshuffle the attention distribution.

When stacked across layers, such all-pass mixing allows UV perturbations in a few tokens to percolate through the entire network, potentially causing large logic drift at the output.

\subsection{UV Perturbations and Hallucination Regimes}

We instantiate the corruption operator $\mathcal{C}$ using several perturbation families that approximate realistic noise:

\begin{itemize}[leftmargin=*]
    \item \textbf{Lexical noise:} synonym substitution, random word insertion/deletion.
    \item \textbf{Syntactic noise:} clause shuffling, distractor clauses that are task-irrelevant.
    \item \textbf{Adversarial noise:} short adversarial prompts or misleading suggestions.
\end{itemize}

These operations are localized and do not alter the ground-truth answer by construction, but they can induce hallucinations in standard LLMs. Our goal is to design architectures that minimize logic drift under these perturbations.

\subsection{Goal}

Formally, given a model class $\mathcal{F}$ and perturbation family $\mathcal{C}$, we seek a model
\begin{equation}
    f^* = \arg\min_{f_\theta \in \mathcal{F}}
    \Big[
        \mathcal{L}_{\text{task}}(f_\theta) + 
        \lambda \, \mathrm{Drift}(f_\theta)
    \Big],
\end{equation}
where $\mathcal{L}_{\text{task}}$ is the standard task loss (e.g., cross-entropy) and $\lambda$ trades off accuracy and robustness. Fermi-Transformers instantiate a particular choice of $\mathcal{F}$ motivated by RG and Fermi statistics.

% =========================================================
% 3. Fermi-Transformer Architecture
% =========================================================
\section{Fermi-Transformer Architecture}

Fermi-Transformers are built by replacing two components of a standard decoder-only Transformer: (i) adding RSMI-Pooling layers to perform semantic coarse-graining, and (ii) substituting softmax with Fermi-Dirac Gapped Attention.

Figure~\ref{fig:architecture} provides an overview of the architecture.

\begin{figure}[t]
    \centering
    % 用 box 画出的占位图
    \fbox{%
        \begin{minipage}[c][0.35\linewidth][c]{0.9\linewidth}
            \centering
            \includegraphics[width=0.95\linewidth]{figs/architecture.png}
        \end{minipage}%
    }
    \caption{
    High-level architecture of the Fermi-Transformer. 
    RSMI-Pooling periodically coarse-grains token blocks into higher-level semantic variables, while Fermi-Dirac Gapped Attention imposes a learnable energy gap on attention weights to suppress UV noise and stabilize the IR reasoning trajectory.
    }
    \label{fig:architecture}
\end{figure}

\subsection{RSMI-Pooling: Semantic Coarse-Graining}

We partition the sequence into non-overlapping blocks $B_j = (x_{(j-1)b+1}, \dots, x_{jb})$ of size $b$. For each block, we learn a coarse-grained representation
\begin{equation}
    h_j = f_\theta(B_j),
\end{equation}
where $f_\theta$ can be a small Transformer or MLP operating on the local tokens.

Following the Real-Space Mutual Information (RSMI) framework~\cite{koch2018mutual}, we choose $f_\theta$ to maximize mutual information between $h_j$ and the neighboring block $B_{j+1}$:
\begin{equation}
    \theta^* = 
    \arg\max_\theta \RSMI(h_j; B_{j+1})
    = \arg\max_\theta I(h_j; B_{j+1}),
\end{equation}
where $I(\cdot;\cdot)$ denotes mutual information. Intuitively, $h_j$ is optimized to retain precisely those features of $B_j$ that are predictive of longer-range context, integrating out locally unpredictable noise.

Since exact MI is intractable, we adopt a variational lower bound using a critic network $T_\phi$:
\begin{equation}
    I(h_j; B_{j+1}) \geq 
    \mathbb{E}_{p(h,B)}[T_\phi(h,B)] -
    \log \mathbb{E}_{p(h)p(B)}[e^{T_\phi(h,B)}],
    \label{eq:mine}
\end{equation}
similar in spirit to MINE-like estimators. In practice, we add an RSMI loss term
\begin{equation}
    \mathcal{L}_{\text{RSMI}} = - \sum_j \hat{I}(h_j; B_{j+1})
\end{equation}
to the training objective with a coefficient $\beta$.

\begin{algorithm}[t]
\caption{RSMI-Pooling Layer (One Step)}
\label{alg:rsmi}
\begin{algorithmic}[1]
\REQUIRE Token embeddings $X \in \mathbb{R}^{T \times d}$, block size $b$, encoder $f_\theta$, critic $T_\phi$
\STATE Partition $X$ into blocks $B_j \in \mathbb{R}^{b \times d}$, $j = 1, \dots, T/b$
\STATE For each $j$, compute coarse variables $h_j = f_\theta(B_j)$
\STATE Sample positive pairs $(h_j, B_{j+1})$ and negative pairs $(h_j, B_{k})$ with $k \neq j+1$
\STATE Estimate mutual information $\hat{I}(h_j; B_{j+1})$ via Eq.~\eqref{eq:mine}
\STATE Form pooled sequence $H = (h_1, \dots, h_{T/b})$
\RETURN Pooled sequence $H$ and RSMI loss $\mathcal{L}_{\text{RSMI}}$
\end{algorithmic}
\end{algorithm}

We insert RSMI-Pooling layers every $k$ Transformer blocks. Subsequent attention layers operate on the shorter sequence $H$, yielding both computational savings and an inductive bias towards long-range semantics.

\subsection{Fermi-Dirac Gapped Attention}

Standard attention weights are computed via softmax:
\begin{equation}
    \alpha_{ij}^{\text{soft}} =
    \frac{\exp(w_{ij})}{\sum_{k} \exp(w_{ik})},
    \quad
    w_{ij} = \frac{q_i^\top k_j}{\sqrt{d_k}}.
\end{equation}
We replace this with a Fermi-Dirac inspired kernel:
\begin{align}
    \tilde{\alpha}_{ij} 
    &= \sigma_F(w_{ij}; \mu_i, T_i)
    = \frac{1}{1 + \exp\!\left(\frac{\mu_i - w_{ij}}{T_i}\right)}, \label{eq:fermi-kernel} \\
    \alpha_{ij} 
    &= \frac{\tilde{\alpha}_{ij}}{\sum_{k} \tilde{\alpha}_{ik}},
\end{align}
where $\mu_i$ is a learnable ``logical Fermi level'' and $T_i$ a temperature, both potentially conditioned on the query representation.

High-scoring interactions with $w_{ij} \gg \mu_i$ saturate to $\tilde{\alpha}_{ij} \approx 1$, while low-scoring interactions with $w_{ij} \ll \mu_i$ are exponentially suppressed. This yields an effective \emph{topological gap} in the attention spectrum.

We define the \emph{Fermi gap} for query position $i$ as
\begin{equation}
    \fermigap{i} =
    \min_{j \in \mathcal{H}_i} w_{ij} -
    \max_{j \in \mathcal{L}_i} w_{ij},
\end{equation}
where $\mathcal{H}_i = \{j : \tilde{\alpha}_{ij} > \tau\}$ denotes the set of ``high-occupancy'' neighbors and $\mathcal{L}_i$ the complement, for some threshold $\tau \in (0.5,1)$.

\begin{proposition}[Gap-Preserving Attention Topology]
\label{prop:gap}
Suppose $\fermigap{i} > 0$ for query position $i$, and consider perturbations $\delta w_{ij}$ with $\|\delta w\|_\infty < \fermigap{i} / 2$. Then the high-occupancy set $\mathcal{H}_i$ induced by Fermi-Dirac Attention remains invariant.
\end{proposition}

\noindent
The proof follows from the fact that the ordering of logits across the gap cannot change under such perturbations, and the Fermi kernel~\eqref{eq:fermi-kernel} is monotone. Proposition~\ref{prop:gap} suggests that as long as the gap remains open, small UV perturbations in logits cannot alter the qualitative topology of attention, providing a form of local rigidity analogous to gapped phases in condensed matter systems~\cite{wen2004quantum}.

\subsection{Theoretical Analysis of Logic Drift}

The gap-preserving property in Proposition~\ref{prop:gap} suggests that, as long as the Fermi gap remains open, the qualitative topology of attention is robust to small perturbations in the logits. We now discuss how this relates to the logic drift metric in Eq.~\eqref{eq:logic-drift} under simplifying assumptions.

Let $A_i \in \mathbb{R}^{T \times T}$ denote the attention matrix for query position $i$, and $V \in \mathbb{R}^{T \times d}$ the value matrix. We write a single-head attention update as
\begin{equation}
    h_i = \sum_{j=1}^{T} \alpha_{ij} v_j,
    \quad
    \alpha_{ij} = \mathrm{Attn}(w_{ij}),
\end{equation}
where $\mathrm{Attn}(\cdot)$ is either softmax or Fermi-Dirac. Consider a perturbed input $\tilde{x}$, inducing perturbed logits $\tilde{w}_{ij} = w_{ij} + \delta w_{ij}$ and value vectors $\tilde{v}_j = v_j + \delta v_j$.

\begin{definition}[Per-Layer Drift]
For a given layer $\ell$, the representation drift at position $i$ is
\begin{equation}
    \Delta h_i^{(\ell)} = 
    h_i^{(\ell)}(x) - h_i^{(\ell)}(\tilde{x}),
\end{equation}
and the global layer drift is
\begin{equation}
    D^{(\ell)} = 
    \Big( \sum_{i=1}^{T_\ell} 
        \|\Delta h_i^{(\ell)}\|_2^2
    \Big)^{1/2}.
\end{equation}
\end{definition}

We consider a simplified setting where:
\begin{itemize}[leftmargin=*]
    \item[(A1)] Value projections $v_j = W^V x_j$ are $L_V$-Lipschitz in the token embeddings.
    \item[(A2)] The readout function $g$ that maps final-layer representations to task outputs is $L_g$-Lipschitz.
    \item[(A3)] For each query $i$ and every layer $\ell$, the Fermi gap $\fermigap{i}^{(\ell)}$ is strictly positive and perturbations satisfy $\|\delta w^{(\ell)}\|_\infty < \fermigap{i}^{(\ell)}/2$.
\end{itemize}

Under (A3), Proposition~\ref{prop:gap} guarantees that the high-occupancy sets $\mathcal{H}_i^{(\ell)}$ are unchanged by the perturbation. As a result, the attention backbone---the subgraph induced by high-occupancy edges---remains fixed.

\begin{theorem}[Logic Drift Bound Under Gapped Attention]
\label{thm:drift-bound}
Under assumptions (A1)--(A3), there exists a constant $C > 0$ depending on the network parameters such that
\begin{equation}
    d\big(f_\theta(x), f_\theta(\tilde{x})\big)
    \;\le\;
    C \, \|E(x) - E(\tilde{x})\|_2,
\end{equation}
where $E(\cdot)$ denotes the token embedding map. In particular, if the perturbation operator $\mathcal{C}$ is such that
$\mathbb{E}\|E(x) - E(\mathcal{C}(x))\|_2 \le \varepsilon$,
then the expected logic drift satisfies
\begin{equation}
    \mathrm{Drift}(f_\theta)
    \;\le\;
    C \, \varepsilon.
\end{equation}
\end{theorem}

\begin{proof}[Proof sketch]
The proof proceeds in three steps. First, under (A3) and Proposition~\ref{prop:gap}, the attention backbone---the set of high-occupancy edges---is invariant under the perturbation. Consequently, the attention update at each layer can be written as a bounded linear operator on the value matrix restricted to a fixed sparsity pattern.

Second, by (A1) and the boundedness of attention weights within each backbone component, the map $E(x) \mapsto h^{(\ell)}(x)$ is Lipschitz for every layer $\ell$, with a constant that depends only on the layer parameters and the gap magnitude. Combining these layer-wise Lipschitz bounds across $L$ layers yields a global Lipschitz constant $L_h$ for the representation map $E(x) \mapsto h^{(L)}(x)$.

Finally, applying (A2) to the readout $g$ and composing with $L_h$ gives a bound of the form
$d(f_\theta(x), f_\theta(\tilde{x})) \le L_g L_h \|E(x) - E(\tilde{x})\|_2$,
from which the statement follows by setting $C = L_g L_h$ and taking expectations.
\end{proof}

Theorem~\ref{thm:drift-bound} should be viewed as a qualitative illustration rather than a tight quantitative guarantee: in practice, real models may violate the simplifying assumptions, and the effective Lipschitz constants can be large. Nevertheless, the result formalizes the intuition that maintaining an open Fermi gap stabilizes the attention topology and helps control logic drift in proportion to the magnitude of token-level perturbations.


\subsection{Training Objective and Complexity}

We train Fermi-Transformers with a combined loss:
\begin{equation}
    \mathcal{L} = 
    \mathcal{L}_{\text{LM}} + 
    \beta \mathcal{L}_{\text{RSMI}},
\end{equation}
where $\mathcal{L}_{\text{LM}}$ is the standard language modeling or task-specific loss, and $\beta$ controls the strength of the RG-inspired regularization. Fermi-Dirac Attention has the same asymptotic complexity as softmax attention ($\mathcal{O}(T^2)$ per layer), while RSMI-Pooling reduces the effective sequence length at higher layers to $T/b$.

% =========================================================
% 4. Experiments
% =========================================================
\section{Experiments}

We evaluate Fermi-Transformers on a suite of synthetic and real-world tasks designed to probe robustness under UV perturbations.

\subsection{Tasks and Datasets}

Dataset statistics (number of examples, average length, and noise rate) are summarized in Table~\ref{tab:datasets}.

\begin{table}[t]
    \centering
    \small
    \caption{Datasets and tasks used in our experiments. Noise rate denotes the fraction of tokens perturbed by the UV corruption operator $\mathcal{C}$. All numbers are placeholders and will be replaced by the final experimental configuration.}
    \label{tab:datasets}
    \begin{tabular}{lcccccc}
        \toprule
        Task & Dataset & Train & Val & Test & Avg Len & Noise \\
        \midrule
        Synthetic reasoning & \TODO{Syn-Reason} & \TODO{50k} & \TODO{5k} & \TODO{5k} & \TODO{64} & \TODO{20\%} \\
        Arithmetic reasoning & \TODO{GSM8K-N} & \TODO{7.5k} & \TODO{0.7k} & \TODO{1k} & \TODO{80} & \TODO{15\%} \\
        Program synthesis & \TODO{HumanEval-N} & \TODO{--} & \TODO{--} & \TODO{164} & \TODO{150} & \TODO{10\%} \\
        Multi-hop QA & \TODO{HotpotQA-N} & \TODO{80k} & \TODO{5k} & \TODO{5k} & \TODO{200} & \TODO{25\%} \\
        Factual QA & \TODO{TruthfulQA-N} & \TODO{--} & \TODO{--} & \TODO{817} & \TODO{30} & \TODO{15\%} \\
        \bottomrule
    \end{tabular}
\end{table}

\paragraph{Synthetic UV-Scrambled Reasoning.}
We construct a symbolic reasoning benchmark where inputs consist of short logical formulas and natural language questions. For each example, we generate noisy variants by corrupting \TODO{$p$\%} of tokens using the perturbations described in Section~2.2. Ground-truth answers remain unchanged.

\paragraph{Noisy Arithmetic and Program Synthesis.}
We adapt existing arithmetic word problem and program synthesis datasets (e.g., \TODO{GSM8K, HumanEval}) by injecting localized paraphrases and irrelevant clauses into the problem statements, simulating noisy scientific environments.

\paragraph{Real-World QA with Adversarial Noise.}
We further evaluate on a subset of multi-hop QA and factuality benchmarks (e.g., \TODO{HotpotQA, TruthfulQA}), where we augment questions with adversarially generated misleading context or irrelevant documents.

\subsection{Models and Training Setup}

We instantiate Fermi-Transformers by modifying a standard decoder-only Transformer:

\begin{itemize}[leftmargin=*]
    \item \textbf{Base architecture.} \TODO{$L$} layers, hidden size \TODO{$d$}, \TODO{$h$} attention heads, feed-forward dimension \TODO{$4d$}.
    \item \textbf{Fermi-Dirac Attention.} $\mu_i$ and $T_i$ are initialized to \TODO{values} and parameterized as small linear projections of the query representation.
    \item \textbf{RSMI-Pooling.} Every \TODO{$k$} layers, we apply RSMI-Pooling with block size $b = \TODO{\text{value}}$ and a shared critic $T_\phi$.
\end{itemize}

We compare against the following baselines:

\begin{itemize}[leftmargin=*]
    \item \textbf{Transformer.} A vanilla Transformer with softmax attention and identical depth and width.
    \item \textbf{Sparse/Long-Range Attention.} A Longformer-style variant with local and global attention~\cite{beltagy2020longformer}.
    \item \textbf{Gated Attention.} A baseline where scalar gates are applied to attention logits without Fermi statistics.
\end{itemize}

All models are trained with the same optimizer, learning rate schedule, batch size, and number of steps. \TODO{Specify optimizer, LR, batch size, training steps, and implementation details.} A summary of key hyperparameters is provided in Table~\ref{tab:hyperparams}.

\begin{table}[t]
    \centering
    \small
    \caption{Key model and training hyperparameters. All values are placeholders and will be replaced by the final configuration.}
    \label{tab:hyperparams}
    \begin{tabular}{lcc}
        \toprule
        Hyperparameter & Value & Notes \\
        \midrule
        Layers $L$ & \TODO{12} & All models \\
        Hidden size $d$ & \TODO{768} & \\
        Heads $h$ & \TODO{12} & \\
        FFN dim & \TODO{3072} & $4d$ \\
        Block size $b$ & \TODO{4} & RSMI-Pooling \\
        Pooling interval $k$ & \TODO{3} & Every $k$ layers \\
        Initial $\mu$ & \TODO{0.0} & Per head \\
        Initial $T$ & \TODO{0.5} & Per head \\
        Optimizer & \TODO{AdamW} & \\
        LR & \TODO{$3\times10^{-4}$} & Cosine decay \\
        Batch size & \TODO{128} & Tokens or sequences \\
        Training steps & \TODO{100k} & \\
        \bottomrule
    \end{tabular}
\end{table}

\subsection{Metrics}

We report:

\begin{itemize}[leftmargin=*]
    \item \textbf{Clean accuracy} on unperturbed test sets.
    \item \textbf{Corrupted accuracy} under UV perturbations.
    \item \textbf{Logic Drift} $\mathrm{Drift}(f_\theta)$ as defined in Eq.~\eqref{eq:logic-drift}, estimated with \TODO{$N$} random perturbations per input.
    \item \textbf{Hallucination rate} on factual QA, measured as the fraction of confidently incorrect answers.
\end{itemize}

\subsection{Main Results}

Table~\ref{tab:main_results} summarizes performance across tasks. We highlight three key observations:

\begin{itemize}[leftmargin=*]
    \item \textbf{Reduced logic drift.} On Synthetic UV-Scrambled Reasoning, Fermi-Transformer reduces logic drift by up to \TODO{xx\%} compared to the vanilla Transformer, and by \TODO{yy\%} on noisy arithmetic tasks.
    \item \textbf{Improved robustness-accuracy trade-off.} On GSM8K-noise and adversarial QA, Fermi-Transformer narrows the clean vs noisy performance gap by \TODO{zz\%} while maintaining comparable clean accuracy.
    \item \textbf{Lower hallucination rate.} On adversarial factual QA, hallucination rate decreases from \TODO{a\%} (Transformer) to \TODO{b\%} (Fermi-Transformer).
\end{itemize}

\begin{table*}[t]
    \centering
    \small
    \caption{
    Main results on clean and noisy benchmarks.
    \textbf{Clean Acc} and \textbf{Noisy Acc} denote accuracy on unperturbed and corrupted inputs, respectively. \textbf{Gap} is the clean--noisy difference (lower is better). \textbf{Drift} is the logic drift metric (lower is better), and \textbf{Halluc.} is hallucination rate on factual QA (lower is better). All numbers are placeholders to be replaced with actual experimental results.
    }
    \label{tab:main_results}
    \begin{tabular}{l l c c c c c}
        \toprule
        Task & Model & Clean Acc $\uparrow$ & Noisy Acc $\uparrow$ & Gap $\downarrow$ & Drift $\downarrow$ & Halluc. $\downarrow$ \\
        \midrule
        Synthetic reasoning &
        Transformer &
        \TODO{85.0} &
        \TODO{62.3} &
        \TODO{22.7} &
        \TODO{0.38} &
        \TODO{--} \\

        & Fermi-Transformer &
        \TODO{86.1} &
        \TODO{75.4} &
        \TODO{10.7} &
        \TODO{0.19} &
        \TODO{--} \\
        \midrule
        GSM8K-noise &
        Transformer &
        \TODO{53.2} &
        \TODO{30.5} &
        \TODO{22.7} &
        \TODO{0.42} &
        \TODO{--} \\

        & Fermi-Transformer &
        \TODO{54.0} &
        \TODO{40.8} &
        \TODO{13.2} &
        \TODO{0.27} &
        \TODO{--} \\
        \midrule
        HotpotQA-noise &
        Transformer &
        \TODO{72.0} &
        \TODO{55.3} &
        \TODO{16.7} &
        \TODO{0.35} &
        \TODO{--} \\

        & Fermi-Transformer &
        \TODO{72.5} &
        \TODO{63.9} &
        \TODO{8.6} &
        \TODO{0.21} &
        \TODO{--} \\
        \midrule
        TruthfulQA-noise &
        Transformer &
        \TODO{--} &
        \TODO{--} &
        \TODO{--} &
        \TODO{0.40} &
        \TODO{35.0} \\

        & Fermi-Transformer &
        \TODO{--} &
        \TODO{--} &
        \TODO{--} &
        \TODO{0.28} &
        \TODO{23.5} \\
        \bottomrule
    \end{tabular}
\end{table*}

\subsection{Ablation Studies}

We perform ablations to disentangle the contributions of each component. Results are shown in Table~\ref{tab:ablation}.

\begin{table}[t]
    \centering
    \small
    \caption{
    Ablation results on \TODO{Synthetic reasoning} and \TODO{GSM8K-noise}. 
    \textsc{No-RSMI} removes RSMI-Pooling,
    \textsc{No-Gap} replaces Fermi-Dirac Attention with softmax.
    All numbers are placeholders.
    }
    \label{tab:ablation}
    \begin{tabular}{l c c c}
        \toprule
        Variant & Task & Noisy Acc $\uparrow$ & Drift $\downarrow$ \\
        \midrule
        Fermi-Transformer & Syn-Reason & \TODO{75.4} & \TODO{0.19} \\
        \textsc{No-RSMI} & Syn-Reason & \TODO{69.8} & \TODO{0.26} \\
        \textsc{No-Gap} & Syn-Reason & \TODO{66.1} & \TODO{0.30} \\
        \midrule
        Fermi-Transformer & GSM8K & \TODO{40.8} & \TODO{0.27} \\
        \textsc{No-RSMI} & GSM8K & \TODO{36.5} & \TODO{0.33} \\
        \textsc{No-Gap} & GSM8K & \TODO{34.0} & \TODO{0.36} \\
        \bottomrule
    \end{tabular}
\end{table}

\subsection{Robustness to Unseen Perturbations}

The previous experiments evaluate robustness under perturbation families that are also present during training (e.g., lexical and syntactic noise). A natural question is whether Fermi-Transformers can generalize to \emph{unseen} UV perturbations, where the corruption mechanism at test time differs from that used during training.

We consider a train--test mismatch setting where models are finetuned with one family of perturbations and evaluated with another:
\begin{itemize}[leftmargin=*]
    \item \textbf{Train on lexical noise, test on syntactic noise.}
    \item \textbf{Train on syntactic noise, test on adversarial prompts.}
    \item \textbf{Train on mixed noise, test on unseen adversarial compositions.}
\end{itemize}

Table~\ref{tab:unseen} reports the noisy accuracy and logic drift on the synthetic reasoning and GSM8K-noise tasks under these mismatched conditions.

\begin{table}[t]
    \centering
    \small
    \caption{
    Robustness to unseen perturbations.
    Models are trained with one perturbation family and evaluated with another. 
    We report noisy accuracy (\%) and logic drift (lower is better) on \TODO{Syn-Reason} and \TODO{GSM8K-noise}. 
    All numbers are placeholders.
    }
    \label{tab:unseen}
    \begin{tabular}{l l c c}
        \toprule
        Train $\rightarrow$ Test & Model & Noisy Acc $\uparrow$ & Drift $\downarrow$ \\
        \midrule
        Lexical $\rightarrow$ Syntactic &
        Transformer &
        \TODO{60.2} &
        \TODO{0.41} \\
        & Fermi-Trans. &
        \TODO{70.5} &
        \TODO{0.27} \\
        \midrule
        Syntactic $\rightarrow$ Adversarial &
        Transformer &
        \TODO{55.8} &
        \TODO{0.45} \\
        & Fermi-Trans. &
        \TODO{64.9} &
        \TODO{0.30} \\
        \midrule
        Mixed $\rightarrow$ Unseen comp. &
        Transformer &
        \TODO{62.0} &
        \TODO{0.38} \\
        & Fermi-Trans. &
        \TODO{69.3} &
        \TODO{0.25} \\
        \bottomrule
    \end{tabular}
\end{table}

We observe that:
\begin{itemize}[leftmargin=*]
    \item Across all train--test mismatches, Fermi-Transformers retain a larger fraction of their robustness, with noisy accuracy improvements of \TODO{5--10} points and logic drift reductions of \TODO{20--30\%} relative to the vanilla Transformer.
    \item The gains are most pronounced in settings where adversarial prompts are introduced only at test time, suggesting that the combination of RSMI-Pooling and Fermi gaps induces an inductive bias that generalizes beyond the specific noise distribution encountered during training.
\end{itemize}

These results indicate that Fermi-Transformers are not merely overfitting to a particular corruption scheme, but instead learn a more noise-tolerant semantic backbone that extends to unseen UV perturbations.

\subsection{Case Study: Attention Topology Under Perturbations}

To provide qualitative intuition, Figure~\ref{fig:case_study} visualizes the attention patterns of a representative example under clean and perturbed prompts, for both the vanilla Transformer and the Fermi-Transformer.

\begin{figure}[t]
    \centering
    % 用 box 画出的占位图
    \fbox{%
        \begin{minipage}[c][0.35\linewidth][c]{0.9\linewidth}
            \centering
            \TODO{Case study attention heatmap placeholder}
        \end{minipage}
    }
    \caption{
    Qualitative case study of attention patterns under UV perturbations.
    Top: vanilla Transformer; attention mass shifts to irrelevant tokens after perturbation, leading to an incorrect answer.
    Bottom: Fermi-Transformer; the high-occupancy attention backbone remains stable thanks to the Fermi gap, and the answer is preserved.
    \TODO{Replace with actual visualization.}
    }
    \label{fig:case_study}
\end{figure}

\subsection{Spectral Analysis}

To probe the connection with RG fixed points, we analyze the singular value spectrum of attention matrices. Let $A_\ell \in \mathbb{R}^{T_\ell \times T_\ell}$ be the average attention matrix at layer $\ell$, with singular values $\sigma_{n}^{(\ell)}$. We find that Fermi-Transformer exhibits a stable power-law decay
\begin{equation}
    \sigma_{n}^{(\ell)} \propto n^{-\alpha}, \quad \alpha \approx \TODO{1.0},
\end{equation}
across layers, whereas the vanilla Transformer shows a more irregular spectrum without a clear scaling regime.

Figure~\ref{fig:spectrum} plots $\sigma_n$ vs.~$n$ in log-log scale for both models, illustrating the emergence of a linear trend in Fermi-Transformers. This empirical observation is consistent with convergence towards a scale-invariant semantic manifold~\cite{zimmern2024large, sethna2001crackling}.

\begin{figure}[t]
    \centering
    % 用 box 画出的占位图
    \fbox{%
        \begin{minipage}[c][0.35\linewidth][c]{0.9\linewidth}
            \centering
            \TODO{Spectral plot placeholder}
        \end{minipage}
    }
    \caption{
    Singular value spectra of attention matrices on a log-log scale.
    Fermi-Transformers exhibit a stable power-law decay with exponent $\alpha \approx \TODO{1.0}$ across layers, indicative of a scale-invariant regime. Vanilla Transformers show a more disordered spectrum.
    \TODO{Replace with final spectrum plot and fitted slopes.}
    }
    \label{fig:spectrum}
\end{figure}

% =========================================================
% 5. Related Work
% =========================================================
\section{Related Work}

\subsection{Statistical Physics and Deep Learning}

Connections between statistical mechanics and neural networks date back to the Hopfield model and its relationship to the Ising system~\cite{hopfield1982neural}. Subsequent work has explored energy landscapes, spin-glass behavior, and self-organized criticality in neural networks and biological systems~\cite{parisi1980order,bak1987self,chialvo2010emergent,biroli2024energy,zimmern2024large}. Our work leverages RG and topological notions to explicitly shape the information flow in Transformers.

\subsection{Renormalization, Information Bottlenecks, and RG Flows}

Mehta and Schwab~\cite{mehta2014exact} formalized an exact mapping between variational RG and deep learning, while Koch-Janusz and Ringel~\cite{koch2018mutual} introduced RSMI as a principled criterion for learning relevant coarse-grained variables. Information-theoretic perspectives such as the information bottleneck principle~\cite{tishby2015deep} and physics-informed learning~\cite{lin2017why,raissi2019physics,merchant2023scaling} have further deepened our understanding of representation learning. Our RSMI-Pooling module directly instantiates an RG-like coarse-graining in token space.

\subsection{Topological and Geometric Deep Learning}

Topological methods have been used to analyze and design neural representations~\cite{carlsson2009topology,bronstein2017geometric,battiloro2024topological,fujimoto2024topological}, while tensor networks and quantum-inspired architectures have offered alternative views on expressivity and scaling~\cite{carleo2017solving,stoudenmire2016supervised,vidal2008class,han2018unsupervised,novikov2015tensor,ali2024tensor}. Our Fermi-Dirac Gapped Attention introduces an explicit gap in the attention spectrum, providing a topological rigidity that is complementary to these approaches.

\subsection{Attention Variants, Long-Context Models, and Robustness}

Numerous variants of attention have been proposed to improve efficiency and handle long contexts~\cite{beltagy2020longformer,katharopoulos2020transformers,dao2022flashattention,gu2023mamba}. Robustness to adversarial prompts and hallucinations is an active area of research, often leveraging data augmentation, decoding strategies, or post-hoc calibration. In contrast, we focus on architectural changes that embed robustness into the information flow itself via RG and Fermi-inspired principles.

% =========================================================
% 6. Discussion and Limitations
% =========================================================
\section{Discussion and Limitations}

\paragraph{Interpretation as a Physical Phase.}
Our results support the view that robust reasoning can be interpreted as a gapped phase of semantic information flow, stabilized by coarse-graining and topological gaps in attention. The emergence of power-law spectra suggests that Fermi-Transformers operate near a scale-invariant regime, potentially related to self-organized criticality in biological neural systems~\cite{chialvo2010emergent}.

\paragraph{Limitations.}
This work has several limitations. First, our experiments are conducted on medium-scale models and benchmarks; it remains to be seen how Fermi-Transformers behave at billion-parameter scales. Second, RSMI estimation introduces additional computational overhead and requires careful tuning of the critic network. Third, imposing strong gaps may conflict with tasks that demand high sensitivity to fine-grained lexical cues.

\paragraph{Future Directions.}
Future work includes scaling Fermi-Transformers to larger models, exploring task-adaptive or layer-wise varying gaps, and extending the framework to multimodal architectures. More broadly, we believe that importing concepts from RG, topology, and non-equilibrium physics into LLM design is a promising direction for achieving robust and interpretable AI systems.

% =========================================================
% 7. Conclusion
% =========================================================
\section{Conclusion}

We introduced Fermi-Transformers, a physics-inspired architecture that combines RSMI-based semantic coarse-graining with Fermi-Dirac Gapped Attention to mitigate UV instabilities in language models. By explicitly integrating out local noise and enforcing a learnable gap in attention, Fermi-Transformers achieve lower logic drift and improved robustness under token-level perturbations, while preserving competitive performance on clean inputs. Spectral analyses further suggest that the model converges towards a scale-invariant regime, aligning with RG intuition.

Our findings indicate that RG and topological ideas can do more than metaphorically describe deep learning: they can inform concrete architectural choices that enhance the robustness of large-scale language models.

\newpage
\bibliographystyle{plainnat}
\bibliography{refs}

\end{document}
